% ======================================================================
% Chapter 11: Renormalization and Symmetry
% (book pages 347-391; PDF pages 370-415)
% ======================================================================
\chapter{Renormalization and Symmetry}

Now that we have determined the general structure of the ultraviolet
divergences of quantum field theories, it would seem natural to continue
investigating the implications of these divergences in Feynman diagram
calculations.  However, we will now put this issue aside until Chapter~12 and
set off in what may seem an unrelated direction.  In Chapter~8 and in
Section~9.3, we noted the formal relation between quantum field theory and
statistical mechanics.  The closest formal analogue of a scalar field theory
was seen to be the continuum description of a ferromagnet or some other system
that allows a second-order phase transition.  This analogy raises the
possibility that in quantum field theory as well it may be possible for the
field to take on a nonzero global value.  As in a magnet, this global field
might have a directional character, and thus violate a symmetry of the
Lagrangian.  In such a case, we say that the field theory has hidden or
\emph{spontaneously broken symmetry}.  We devote this chapter to an analysis
of this mechanism of symmetry violation.

Spontaneously broken symmetry is a central concept in the study of quantum
field theory, for two reasons.  First, it plays a major role in the
applications of quantum field theory to Nature.  In this book, we will see two
very different examples of such applications: Chapter~13 will apply the theory
of hidden symmetry to statistical mechanics, specifically to the behavior of
thermodynamic variables near second-order phase transitions.  Later, in
Chapter~20, we will see that hidden symmetry is an essential ingredient in the
theory of the weak interactions.  Spontaneous symmetry breaking also finds
applications in the theory of the strong interactions, and in the search for
unified models of fundamental physics.

But spontaneous symmetry breaking is also interesting from a theoretical point
of view.  Quantum field theories with spontaneously broken symmetry contain
ultraviolet divergences.  Thus, it is natural to ask whether these divergences
are constrained by the underlying symmetry of the theory.  The answer to this
question, first presented by Benjamin Lee,%
\footnote{A beautiful summary of Lee's analysis is given in his lecture note
volume: B.~Lee, \emph{Chiral Dynamics} (Gordon and Breach, New York, 1972).}
will give us further insights into the nature of ultraviolet divergences and
the meaning of renormalization.

\section{Spontaneous Symmetry Breaking}

We begin with an analysis of spontaneous symmetry breaking in classical field
theory.  Consider first the familiar $\phi^4$ theory Lagrangian,
\begin{equation}
  \mathcal{L} = \tfrac{1}{2}(\partial_\mu\phi)^2 - \tfrac{1}{2}m^2\phi^2
    - \frac{\lambda}{4!}\phi^4 ,
\end{equation}
but with $m^2$ replaced by a negative parameter, $-\mu^2$:
\begin{equation}
  \mathcal{L} = \tfrac{1}{2}(\partial_\mu\phi)^2 + \tfrac{1}{2}\mu^2\phi^2
    - \frac{\lambda}{4!}\phi^4 .
  \label{eq:11.1}
\end{equation}
This Lagrangian has a discrete symmetry: It is invariant under the operation
$\phi\to-\phi$.  The corresponding Hamiltonian is
\begin{equation}
  H = \int d^3x\ \Bigl[ \tfrac{1}{2}\pi^2 + \tfrac{1}{2}(\nabla\phi)^2
    + V(\phi) \Bigr] .
\end{equation}
The minimum-energy classical configuration is a uniform field
$\phi(x) = \phi_0$, with $\phi_0$ chosen to minimize the potential
\begin{equation}
  V(\phi) = -\frac{1}{2}\mu^2\phi^2 + \frac{\lambda}{4!}\phi^4
\end{equation}
(see Fig.~\ref{fig:11.1}).  This potential has two minima, given by
\begin{equation}
  \phi_0 = \pm v = \pm \sqrt{\frac{6\mu^2}{\lambda}} .
  \label{eq:11.2}
\end{equation}
The constant $v$ is called the \emph{vacuum expectation value} of $\phi$.

\begin{figure}[t]
  \centering
  \includegraphics[width=0.45\textwidth]{fig11.1.png}
  \caption{Potential for spontaneous symmetry breaking in the discrete case.}
  \label{fig:11.1}
\end{figure}

To interpret this theory, suppose that the system is near one of the minima
(say the positive one).  Then it is convenient to define
\begin{equation}
  \phi(x) = v + \sigma(x),
  \label{eq:11.3}
\end{equation}
and rewrite $\mathcal{L}$ in terms of $\sigma(x)$.  Plugging \eqref{eq:11.3}
into \eqref{eq:11.1}, we find that the term linear in $\sigma$ vanishes (as it
must, since the minimum of the potential is at $\sigma = 0$).  Dropping the
constant term as well, we obtain the Lagrangian
\begin{equation}
  \mathcal{L} = \tfrac{1}{2}(\partial_\mu\sigma)^2
    - \tfrac{1}{2}(2\mu^2)\sigma^2 - \sqrt{\frac{\lambda}{6}}\,\mu\,\sigma^3
    - \frac{\lambda}{4!}\sigma^4 .
  \label{eq:11.4}
\end{equation}
This Lagrangian describes a simple scalar field of mass $\sqrt{2}\mu$ with
$\sigma^3$ and $\sigma^4$ interactions.  The symmetry $\phi\to-\phi$ is no
longer apparent; its only manifestation is in the relations among the three
coefficients in \eqref{eq:11.4}, which depend in a special way on only two
parameters.  This is the simplest example of a spontaneously broken symmetry.

\subsection{The Linear Sigma Model}

A more interesting theory arises when the broken symmetry is continuous,
rather than discrete.  The most important example is a generalization of the
preceding theory called the \emph{linear sigma model}, which we considered
briefly in Problem~4.3.  We will study this model in detail throughout this
chapter.

The Lagrangian of the linear sigma model involves a set of $N$ real scalar
fields $\phi^i(x)$,
\begin{equation}
  \mathcal{L} = \tfrac{1}{2}(\partial_\mu\phi^i)^2 + \tfrac{1}{2}\mu^2(\phi^i)^2
    - \frac{\lambda}{4}\bigl[ (\phi^i)^2 \bigr]^2 ,
  \label{eq:11.5}
\end{equation}
with an implicit sum over $i$ in each factor $((\phi^i)^2)$.  Note that we
have rescaled the coupling $\lambda$ from the $\phi^4$ theory Lagrangian to
remove the awkward factors of 6 in the analysis above.  The Lagrangian
\eqref{eq:11.5} is invariant under the symmetry
\begin{equation}
  \phi^i \to R^{ij}\phi^j ,
  \label{eq:11.6}
\end{equation}
for any $N\times N$ orthogonal matrix $R$.  The group of transformations
\eqref{eq:11.6} is just the rotation group in $N$ dimensions, also called the
$N$-dimensional orthogonal group or simply $O(N)$.

Again the lowest-energy classical configuration is a constant field
$\phi^i_0$, whose value is chosen to minimize the potential
\begin{equation}
  V(\phi) = -\frac{1}{2}\mu^2(\phi^i)^2 + \frac{\lambda}{4}\bigl[
    (\phi^i)^2 \bigr]^2
\end{equation}
(see Fig.~\ref{fig:11.2}).  This potential is minimized for any $\phi^i_0$
that satisfies
\begin{equation}
  (\phi^i_0)^2 = \frac{\mu^2}{\lambda} .
\end{equation}
This condition determines only the length of the vector $\phi_0$; its
direction is arbitrary.  It is conventional to choose coordinates so that
$\phi^i_0$ points in the $N$th direction:
\begin{equation}
  \phi^i_0 = (0, 0, \ldots, 0, v),
  \qquad\text{where } v = \sqrt{\frac{\mu^2}{\lambda}} .
  \label{eq:11.7}
\end{equation}
We can now define a set of shifted fields by writing
\begin{equation}
  \phi^i(x) = \bigl( \pi^k(x), v + \sigma(x) \bigr),
  \qquad k = 1, \ldots, N-1 .
  \label{eq:11.8}
\end{equation}
(The notation, as in Problem~4.3, comes from the application of this
formalism to pions in the case $N = 4$.)

\begin{figure}[t]
  \centering
  \includegraphics[width=0.45\textwidth]{fig11.2.png}
  \caption{Potential for spontaneous breaking of a continuous $O(N)$
    symmetry, drawn for the case $N = 2$.  Oscillations along the trough in
    the potential correspond to the massless $\pi$ fields.}
  \label{fig:11.2}
\end{figure}

It is now straightforward to rewrite the Lagrangian \eqref{eq:11.5} in terms
of the $\pi$ and $\sigma$ fields.  The result is
\begin{equation}
  \mathcal{L} = \tfrac{1}{2}(\partial_\mu\pi^k)^2 + \tfrac{1}{2}
    (\partial_\mu\sigma)^2 - \frac{1}{2}(2\mu^2)\sigma^2
    - \sqrt{\lambda}\,\mu\,\sigma^3 - \sqrt{\lambda}\,\mu\,(\pi^k)^2\sigma
    - \frac{\lambda}{4}\sigma^4 - \frac{\lambda}{2}(\pi^k)^2\sigma^2
    - \frac{\lambda}{4}\bigl[ (\pi^k)^2 \bigr]^2 .
  \label{eq:11.9}
\end{equation}
We obtain a massive $\sigma$ field just as in \eqref{eq:11.4}, and also a set
of $N-1$ massless $\pi$ fields.  The original $O(N)$ symmetry is hidden,
leaving only the subgroup $O(N-1)$, which rotates the $\pi$ fields among
themselves.  Referring to Fig.~\ref{fig:11.2}, we note that the massive
$\sigma$ field describes oscillations of $\phi^i$ in the radial direction, in
which the potential has a nonvanishing second derivative.  The massless $\pi$
fields describe oscillations of $\phi^i$ in the tangential directions, along
the trough of the potential.  The trough is an $(N-1)$-dimensional surface,
and all $N-1$ directions are equivalent, reflecting the unbroken $O(N-1)$
symmetry.

\subsection{Goldstone's Theorem}

The appearance of massless particles when a continuous symmetry is
spontaneously broken is a general result, known as \emph{Goldstone's theorem}.
To state the theorem precisely, we must count the number of linearly
independent continuous symmetry transformations.  In the linear sigma model,
there are no continuous symmetries for $N = 1$, while for $N = 2$ there is a
single direction of rotation.  A rotation in $N$ dimensions can be in any one
of $N(N-1)/2$ planes, so the $O(N)$-symmetric theory has $N(N-1)/2$
continuous symmetries.  After spontaneous symmetry breaking there are
$(N-1)(N-2)/2$ remaining symmetries, corresponding to rotations of the
$(N-1)$ $\pi$ fields.  The number of broken symmetries is the difference,
$N-1$.

Goldstone's theorem states that for every spontaneously broken continuous
symmetry, the theory must contain a massless particle.%
\footnote{J.~Goldstone, \emph{Nuovo Cim.} \textbf{19}, 154 (1961).  An
instructive four-page paper by J.~Goldstone, A.~Salam, and S.~Weinberg,
\emph{Phys.\ Rev.} \textbf{127}, 965 (1962), gives three different proofs of
the theorem.}  We have just seen that this theorem holds in the linear sigma
model, at least at the classical level.  The massless fields that arise through
spontaneous symmetry breaking are called \emph{Goldstone bosons}.  Many light
bosons seen in physics, such as the pions, may be interpreted (at least
approximately) as Goldstone bosons.  We conclude this section with a general
proof of Goldstone's theorem for classical scalar field theories.  The rest of
this chapter is devoted to the quantum-mechanical analysis of theories with
hidden symmetry.  By the end of the chapter we will see that Goldstone bosons
cannot acquire mass from any order of quantum corrections.

Consider, then, a theory involving several fields $\phi^a(x)$, with a
Lagrangian of the form
\begin{equation}
  \mathcal{L} = (\text{terms with derivatives}) - V(\phi) .
  \label{eq:11.10}
\end{equation}
Let $\phi_0$ be a constant field that minimizes $V$, so that
\begin{equation}
  \frac{\partial V}{\partial\phi^a}\Big|_{\phi_0} = 0 .
\end{equation}
Expanding $V$ about this minimum, we find
\begin{equation}
  V(\phi) = V(\phi_0) + \frac{1}{2}(\phi-\phi_0)^a\,
    \frac{\partial^2 V}{\partial\phi^a\partial\phi^b}\Big|_{\phi_0}\,
    (\phi-\phi_0)^b + \cdots .
\end{equation}
The coefficient of the quadratic term,
\begin{equation}
  M^2_{ab} = \frac{\partial^2 V}{\partial\phi^a\partial\phi^b}\Big|_{\phi_0},
  \label{eq:11.11}
\end{equation}
is a symmetric matrix whose eigenvalues give the masses of the fields.  These
eigenvalues cannot be negative, since $\phi_0$ is a minimum.  To prove
Goldstone's theorem, we must show that every continuous symmetry of the
Lagrangian \eqref{eq:11.10} that is not a symmetry of $\phi_0$ gives rise to a
zero eigenvalue of this mass matrix.

A general continuous symmetry transformation has the form
\begin{equation}
  \phi^a \to \phi^a + \alpha\Delta^a(\phi),
  \label{eq:11.12}
\end{equation}
where $\alpha$ is an infinitesimal parameter and $\Delta^a$ is some function of
all the $\phi$'s.  Specialize to constant fields; then the derivative terms in
$\mathcal{L}$ vanish and the potential alone must be invariant under
\eqref{eq:11.12}.  This condition can be written
\begin{equation}
  V(\phi^a) = V\bigl( \phi^a + \alpha\Delta^a(\phi) \bigr)
  \qquad\text{or}\qquad
  \Delta^a(\phi)\frac{\partial V}{\partial\phi^a} = 0 .
\end{equation}
Now differentiate with respect to $\phi^b$, and set $\phi = \phi_0$:
\begin{equation}
  0 = \Bigl( \frac{\partial\Delta^a}{\partial\phi^b} \frac{\partial V}
    {\partial\phi^a} + \Delta^a \frac{\partial^2 V}{\partial\phi^a
      \partial\phi^b} \Bigr)\Big|_{\phi_0}
  = \Delta^a(\phi_0)\,M^2_{ab} .
  \label{eq:11.13}
\end{equation}
The first term vanishes since $\phi_0$ is a minimum of $V$, so the second term
must also vanish.  If the transformation leaves $\phi_0$ unchanged (i.e., if
the symmetry is respected by the ground state), then $\Delta^a(\phi_0) = 0$
and this relation is trivial.  A spontaneously broken symmetry is precisely
one for which $\Delta^a(\phi_0) \neq 0$; in this case $\Delta^a(\phi_0)$ is our
desired vector with eigenvalue zero, so Goldstone's theorem is proved.

\section{Renormalization and Symmetry: An Explicit Example}

Now let us investigate the quantum mechanics of a theory with spontaneously
broken symmetry.  Again we will use as our example the linear sigma model.
The Lagrangian of this theory, written in terms of shifted fields, is given in
Eq.~\eqref{eq:11.9}.  From this expression, we can read off the Feynman
rules; these are shown in Fig.~\ref{fig:11.3}.

\begin{figure}[t]
  \centering
  \includegraphics[width=0.6\textwidth]{fig11.3.png}
  \caption{Feynman rules for the linear sigma model in terms of shifted
    fields.}
  \label{fig:11.3}
\end{figure}

Using these Feynman rules, we can compute tree-level amplitudes without
difficulty.  Diagrams with loops, however, will often diverge.  For the
amplitude with $N_e$ external legs, the superficial degree of divergence is
\begin{equation}
  D = 4 - N_e ,
\end{equation}
just as in the discussion of $\phi^4$ theory in Section~10.2.  (Diagrams
containing a three-point vertex will be less divergent than this expression
indicates, because this vertex has a coefficient with dimensions of mass.)
However, the symmetry constraints on the amplitudes are much weaker than in
that earlier discussion: The theory has three couplings ($\mu$, $\lambda$, and
$v$) whose relation is fixed only at tree level, and the divergences of the
theory may or may not be consistent with the $O(N)$ symmetry.

The divergent amplitudes are shown in Fig.~\ref{fig:11.4}.  There are
divergences in the vacuum energy, the one-point $\sigma$ amplitude (tadpole),
the two-point functions, and the four-point functions.  If the theory is to be
renormalizable, all of these divergences must be canceled by the counterterms
of the form
\begin{equation}
  \delta\mathcal{L} = \tfrac{1}{2}\delta_Z(\partial_\mu\phi^i)^2
    + \tfrac{1}{2}\delta_\mu\phi^i\phi^i
    + \delta_\lambda\bigl( \phi^i\phi^i \bigr)^2
    + \delta_v\phi^N ,
  \label{eq:11.15}
\end{equation}
which contain only three adjustable parameters, $\delta_Z$, $\delta_\mu$, and
$\delta_\lambda$, together with the tadpole counterterm $\delta_v$.  It is not
at all obvious that these three counterterms suffice to cancel all of the
infinities in Fig.~\ref{fig:11.4}.  This is the question posed at the
beginning of this chapter: Are the divergences of a theory with spontaneously
broken symmetry constrained by the underlying $O(N)$ symmetry?

\begin{figure}[t]
  \centering
  \includegraphics[width=0.6\textwidth]{fig11.4.png}
  \caption{Divergent amplitudes in the linear sigma model.}
  \label{fig:11.4}
\end{figure}

\begin{figure}[t]
  \centering
  \includegraphics[width=0.65\textwidth]{fig11.5.png}
  \caption{Feynman rules for counterterm vertices in the linear sigma model.}
  \label{fig:11.5}
\end{figure}

Fortunately, the miracle does occur.  We will see below that the
counterterms of \eqref{eq:11.15}, even though they contain only three
adjustable parameters, are indeed sufficient to cancel all the infinities that
occur in this theory.  In this section we will demonstrate this cancellation
explicitly at the one-loop level.  The rest of this chapter is devoted to a
more general discussion of these issues.

\subsection{Renormalization Conditions}

In the discussion to follow, we will keep track of only the divergent parts of
Feynman diagrams.  However, it will be useful to keep in mind a set of
renormalization conditions that could, in principle, be used to determine also
the finite parts of the counterterms.  Since the counterterms contain three
adjustable parameters, we need three conditions.  We could take these to be
the conditions \eqref{eq:10.19} (implemented according to \eqref{eq:10.28}),
specifying the physical mass $m$ of the $\sigma$ field, its field strength, and
the scattering amplitude at threshold.  However, it is technically easier to
replace one of these conditions with a constraint on the one-point amplitude
for $\sigma$ (the sum of tadpole diagrams):
\begin{equation}
  \text{(sum of all tadpole diagrams)} = 0 .
  \label{eq:11.16}
\end{equation}
In QED the tadpole diagrams automatically vanish, as we saw in Eq.~\eqref{eq:10.5}.
In the linear sigma model, however, no symmetry forbids the appearance of a
nonvanishing one-$\sigma$ amplitude.  This amplitude produces a vacuum
expectation value of $\sigma$ and so, since $\phi^N = v + \sigma$, shifts the
vacuum expectation value of $\phi$.  Such a shift is quite acceptable, as long
as it is finite after counterterms are properly added into the computation of
the amplitude.  However, it will simplify the bookkeeping to set up our
conventions so that the relation \eqref{eq:11.16} is satisfied to all orders
in perturbation theory.

We will define $\lambda$, as in Eq.~\eqref{eq:10.19}, as the scattering
amplitude at threshold.  Then Eq.~\eqref{eq:11.16} defines the parameter
$\mu$, so the mass $m$ of the $\sigma$ field will differ from the result of the
classical equations $m^2 = 2\mu^2 = 2\lambda v^2$ by terms of order
$\mathcal{O}(\lambda\mu^2)$.  If indeed we can remove the divergences from the
theory by adjusting three counterterms, these corrections will be finite and
constitute a prediction of the quantum field theory.

To summarize, we will use the following renormalization conditions:
\begin{equation}
  \begin{aligned}
    \text{(tadpole)} &= 0;\qquad
    \frac{d\Sigma}{dp^2}\Big|_{p^2=m^2} = 0;\\
    \text{(two-point function)} &\to \frac{i}{p^2-m^2}
      \text{ at } p^2 = m^2 ;\\
    \text{(four-point amplitude)} &= -i\lambda \text{ at } s = 4m^2,\
      t = u = 0 .
  \end{aligned}
  \label{eq:11.17}
\end{equation}
In the last condition, the circle is the amputated four-point amplitude.
Note that the last two conditions depend on the physical mass $m$ of the
$\sigma$ particle.  We must now show that these three conditions suffice to
make all of the one-loop amplitudes of the linear sigma model finite.

\subsection{The Vertex Counterterm}

We begin by determining the counterterm $\delta_\lambda$ by computing the
$4\sigma$ amplitude.  The tree-level term comes from the $4\sigma$ vertex, and
is just such as to satisfy \eqref{eq:11.17}.  The one-loop contribution to
this amplitude is the sum of diagrams:
\begin{equation}
  \text{(one $\sigma$-loop)} + \text{(one $\pi$-loop)} + \text{(crossed
    diagrams)} .
  \label{eq:11.18}
\end{equation}
According to \eqref{eq:11.17}, we must adjust $\delta_\lambda$ so that this
sum of diagrams vanishes at threshold.  In this calculation, we will only keep
track of the ultraviolet divergences.  This greatly simplifies the analysis,
because most of the diagrams in \eqref{eq:11.18} are finite.  All the diagrams
with loops made of three or more propagators are finite, since they have at
least six powers of the loop momentum in the denominator; for example,
\begin{equation}
  \int \frac{d^4k}{(2\pi)^4}\ \frac{1}{k^2}\frac{1}{k^2}\frac{1}{k^2}
  \quad \text{is finite} .
\end{equation}
Alternatively, we can see that this diagram is finite in the following way:
Each three-point vertex carries a factor of $\mu$, which has dimensions of
mass.  According to the dimensional analysis argument of Section~10.1, each
such factor lowers the degree of divergence of a diagram by 1.  Since the
$4\sigma$ amplitude already has $D = 0$, any diagram containing a three-point
vertex must be finite.

We are left with the first two diagrams of \eqref{eq:11.18} and the four
diagrams related to these by crossing.  Let us evaluate the first diagram
using dimensional regularization:
\begin{equation}
  \begin{aligned}
    \text{(one $\sigma$-loop)} &= \frac{1}{2}(-6i\lambda)^2\ \int
      \frac{d^dk}{(2\pi)^d}\ \frac{i}{k^2-2\mu^2}\ \frac{i}{(k+p)^2-2\mu^2}\\
    &= 18i\lambda^2\ \int_0^1 dx\ \frac{\Gamma(2-d/2)}{(4\pi)^{d/2}}\
      \frac{1}{\Delta^{2-d/2}}\\
    &= 18i\lambda^2\ \frac{\Gamma(2-d/2)}{(4\pi)^2}\ + (\text{finite terms}).
  \end{aligned}
  \label{eq:11.19}
\end{equation}
Here $\Delta$ is a function of $p$ and $\mu$, whose exact form does not
concern us.  Since our objective is only to demonstrate the cancellation of
the divergences, we will neglect finite terms here and throughout the rest of
this section.  The second diagram of \eqref{eq:11.18} (with $\pi$'s instead
of $\sigma$'s for the internal lines) is identical, except that each vertex
factor is changed from $-6i\lambda$ to $-2i\lambda\delta^{ij}$.  (Roman
indices $i, j, \ldots$ run from 1 to $N-1$.)  We therefore have
\begin{equation}
  \text{(one $\pi$-loop)} = -2i\lambda^2(N-1)\ \frac{\Gamma(2-d/2)}
    {(4\pi)^2}\ + (\text{finite terms}) .
  \label{eq:11.20}
\end{equation}
Since the infinite part of each of these diagrams is simply a
momentum-independent constant, the infinite parts of the corresponding $t$-
and $u$-channel diagrams must be identical.  Therefore the infinite part of
the $4\sigma$ vertex is just three times the sum of \eqref{eq:11.19} and
\eqref{eq:11.20}:
\begin{equation}
  \text{(infinite part of } 4\sigma \text{ vertex)} \sim 3i\lambda^2(N+8)\
    \frac{\Gamma(2-d/2)}{(4\pi)^2} .
  \label{eq:11.21}
\end{equation}
(In this section we use the $\sim$ symbol to indicate equality up to omitted
finite corrections.)  Applying the third condition of \eqref{eq:11.17}, we
find that the counterterm $\delta_\lambda$ is given by
\begin{equation}
  \delta_\lambda \sim \lambda^2(N+8)\ \frac{\Gamma(2-d/2)}{(4\pi)^2} .
  \label{eq:11.22}
\end{equation}

Once we have determined the value of $\delta_\lambda$, we have fixed the
counterterms for the two other four-point amplitudes.  Are these amplitudes
also made finite?  Consider the amplitude with two $\sigma$'s and two
$\pi$'s.  This receives one-loop corrections from
\begin{equation}
  \text{(one $\sigma$-loop)} + \text{(two $\pi$-loops)} + \text{(one
    $\sigma\pi$ loop)} ,
  \label{eq:11.23}
\end{equation}
and from several diagrams with three-point vertices which, as argued earlier,
are manifestly finite.  Each of the diagrams in \eqref{eq:11.23} contains a
loop integral analogous to that in \eqref{eq:11.19}, whose infinite part is
always $-i\Gamma(2-d/2)/(4\pi)^2$.  The only differences are in the vertices
and symmetry factors.  For example, the infinite part of the first diagram of
\eqref{eq:11.23} is
\begin{equation}
  \sim i\cdot(-6i\lambda)(-2i\lambda\delta^{ij})\cdot
    \frac{1}{2}\ \frac{\Gamma(2-d/2)}{(4\pi)^2}
  = -6i\lambda^2\delta^{ij}\ \frac{\Gamma(2-d/2)}{(4\pi)^2} .
\end{equation}
The second diagram is a bit more complicated:
\begin{equation}
  \begin{aligned}
    &\sim \frac{1}{2}(-2i\lambda\delta^{kl})(-2i\lambda(\delta^{ij}\delta^{kl}
      + \delta^{ik}\delta^{jl} + \delta^{il}\delta^{jk}))\cdot
        \frac{\Gamma(2-d/2)}{(4\pi)^2}\\
    &= 2i\lambda^2(N+1)\delta^{ij}\ \frac{\Gamma(2-d/2)}{(4\pi)^2} .
  \end{aligned}
\end{equation}
In the third diagram there is no symmetry factor:
\begin{equation}
  \sim (-2i\lambda\delta^{ik})(-2i\lambda\delta^{kj})\cdot
    \frac{\Gamma(2-d/2)}{(4\pi)^2}
  = -4i\lambda^2\delta^{ij}\ \frac{\Gamma(2-d/2)}{(4\pi)^2} .
\end{equation}
The fourth diagram of \eqref{eq:11.23} gives an identical expression, since it
is the same as the third but with $i$ and $j$ interchanged.  The sum of the
four diagrams therefore gives, for the infinite part of the
$\sigma\sigma\pi\pi$ vertex,
\begin{equation}
  \sim 2i\lambda^2\bigl( 3(N+8) \bigr)\ \frac{\Gamma(2-d/2)}{(4\pi)^2}
  \ \cdots ,
  \label{eq:11.24}
\end{equation}
This divergent term is indeed canceled by the $\sigma\sigma\pi\pi$
counterterm, with the value of $\delta_\lambda$ given in \eqref{eq:11.22}.

The remaining four-point amplitude has four external $\pi$ fields.  The
divergent one-loop diagrams are:
\begin{equation}
  \text{(two $\pi$-loops)} + \text{(two $\sigma$-exchange diagrams)}
    + \text{(crossed)} .
\end{equation}
These diagrams all have the same familiar form.  The first is
\begin{equation}
  \frac{1}{2}(-2i\lambda\delta^{ij})(-2i\lambda\delta^{ij})\cdot
    \frac{\Gamma(2-d/2)}{(4\pi)^2}
  = -2i\lambda^2\delta^{ij}\ \frac{\Gamma(2-d/2)}{(4\pi)^2} ,
  \label{eq:11.25}
\end{equation}
and similarly for the others.  Summing all the divergent diagrams and adding
the counterterm, one finds that the $4\pi$ amplitude is rendered finite by
the same counterterm $\delta_\lambda$.

\subsection{The Tadpole and the $\sigma$ Mass}

Let us next consider the tadpole diagram and the two-point function.  The
tadpole (one-point) amplitude receives contributions from
\begin{equation}
  \text{(one $\sigma$-loop)} + \text{(one $\pi$-loop)} + \text{(counterterm)} .
  \label{eq:11.30}
\end{equation}
It is convenient to introduce a small mass $\xi$ for the $\pi$ field as an
infrared regulator.  The first diagram is
\begin{equation}
  \text{(one $\sigma$-loop)} \sim \frac{1}{2}(-6i\lambda v)\ \int
    \frac{d^dk}{(2\pi)^d}\ \frac{i}{k^2-2\mu^2}
  \sim -3i\lambda v\ \frac{\Gamma(1-d/2)}{(4\pi)^{d/2}}\ (2\mu^2)^{d/2-1} .
  \label{eq:11.31}
\end{equation}
The second diagram is
\begin{equation}
  \text{(one $\pi$-loop)} \sim \frac{1}{2}(-2i\lambda v\delta^{ii})\ \int
    \frac{d^dk}{(2\pi)^d}\ \frac{i}{k^2-\xi^2}
  \sim -i\lambda v(N-1)\ \frac{\Gamma(1-d/2)}{(4\pi)^{d/2}}\ \xi^{d-2} .
  \label{eq:11.32}
\end{equation}
Notice that, for $d > 2$, the diagram vanishes in the limit as $\xi\to 0$;
however, it has a pole at $d = 2$.  Despite these strange features, we can add
\eqref{eq:11.32} to \eqref{eq:11.31} and impose the condition that the
tadpole diagrams be canceled by the counterterm from Fig.~\ref{fig:11.5}.
This condition gives
\begin{equation}
  \delta_v + (\text{loop diagrams}) = -\frac{\lambda v}{2}\ \frac{\Gamma(1-d/2)}
    {(4\pi)^{d/2}}\ \Bigl( (2\mu^2)^{1-d/2} + (N-1)\xi^{2-d/2} \Bigr) .
  \label{eq:11.33}
\end{equation}

Now consider the $2\sigma$ amplitude.  The one-particle-irreducible amplitude
receives contributions from four one-loop diagrams and a counterterm:
\begin{equation}
  \text{(four one-loop diagrams)} + \text{(counterterm)} .
  \label{eq:11.34}
\end{equation}
It is convenient to write the counterterm vertex as
\begin{equation}
  -i\bigl( 2v^2\delta_\lambda \bigr) - i\bigl( \delta_\mu + v^2\delta_\lambda
    \bigr) + ip^2\delta_Z .
  \label{eq:11.35}
\end{equation}
In a general renormalization scheme, the $\sigma$ mass will also be shifted by
the tadpole diagrams (and their counterterm):
\begin{equation}
  \text{(tadpole corrections to the mass)} .
  \label{eq:11.36}
\end{equation}
However, the first renormalization condition in \eqref{eq:11.17} forces these
diagrams to cancel precisely.  This is an example of the special simplicity of
this renormalization condition.

The first two diagrams of \eqref{eq:11.34} are again manifestations of the
generic four-point diagram, now with two external legs replaced by the vacuum
expectation value of $\phi$.  In analogy with the preceding calculations, we
find for the first diagram
\begin{equation}
  \sim 18i\lambda^2v^2\ \frac{\Gamma(2-d/2)}{(4\pi)^2} ,
\end{equation}
and for the second diagram
\begin{equation}
  \sim 2i\lambda^2v^2(N-1)\ \frac{\Gamma(2-d/2)}{(4\pi)^2} .
\end{equation}
Using \eqref{eq:11.22}, we see that these two contributions are canceled by
the first term of \eqref{eq:11.35}.  The third and fourth diagrams of
\eqref{eq:11.34} contain precisely the same integrals as the tadpole diagrams
of \eqref{eq:11.30}.  Relation \eqref{eq:11.33} implies that they are
canceled by the second term in \eqref{eq:11.35}.  Notice that there is no
divergent term proportional to $p^2$ in any of the one-loop diagrams of
\eqref{eq:11.34}.  Thus the renormalization constant $\delta_Z$ is zero at the
one-loop level, just as in ordinary $\phi^4$ theory.

There remains only one potentially divergent amplitude---the $\pi\pi\pi\pi$
amplitude:
\begin{equation}
  \text{(three one-loop diagrams)} + \text{(counterterm)} .
  \label{eq:11.37}
\end{equation}
In analogy with \eqref{eq:11.31}, the first diagram is
\begin{equation}
  \sim \frac{1}{2}(-2i\lambda\delta^{ij})\ \frac{\Gamma(1-d/2)}
    {(4\pi)^{d/2}}\ \cdots .
\end{equation}
The second diagram is quite similar.  As in \eqref{eq:11.32}, it is useful to
introduce a small pion mass as an infrared regulator.  The third diagram is
given by
\begin{equation}
  \sim (-2i\lambda v\delta^{ik})(-2i\lambda v\delta^{kj})\ \int
    \frac{d^dk}{(2\pi)^d}\ \frac{1}{k^2-2\mu^2}\ \frac{1}{(k+p)^2-\xi^2}
  = 4i\lambda^2v^2\delta^{ij}\ \frac{\Gamma(2-d/2)}{(4\pi)^{d/2}}\
    \int_0^1 dx\ \Bigl( \frac{1}{2\mu^2x + \xi^2(1-x) - p^2x(1-x)}
      \Bigr)^{2-d/2} .
  \label{eq:11.38}
\end{equation}
The divergent part of this expression is independent of $p$, so to check the
cancellation of the divergence, it suffices to set $p = 0$.  It will be
instructive to compute the complete amplitude at $p = 0$, including the
finite terms.  Adding the three loop diagrams and the counterterm, whose
value is given by \eqref{eq:11.33}, we find that the $\pi\pi$ amplitude at
$p = 0$---which is precisely the mass shift $\delta m_\pi^2$ of the $\pi$
field---is not only finite, but vanishes completely:
\begin{equation}
  \delta m_\pi^2 \Big|_{p=0} = 0 .
  \label{eq:11.39}
\end{equation}
This result is very attractive.  The $\pi\pi$ amplitude, at $p = 0$, is
precisely the mass shift $\delta m_\pi^2$ of the $\pi$ field.  We already knew
that the $\pi$ particles are massless at tree level---they are the $N-1$
massless bosons required by Goldstone's theorem.  We have now verified that
these bosons remain massless at the one-loop level in the linear sigma model;
in other words, the first quantum corrections to the linear sigma model also
respect Goldstone's theorem.  At the end of this chapter, we will give a
general argument that Goldstone's theorem is satisfied to all orders in
perturbation theory.

\section{The Effective Action}

In the first section of this chapter, we analyzed spontaneous symmetry
breaking in classical field theory.  That analysis was geometrical: We found
the vacuum state by finding the deepest well in a potential surface, and we
proved Goldstone's theorem by showing that symmetry required the presence of
a line of degenerate minima at the bottom of the well.  But this geometrical
picture was lost, or at least disguised, in the one-loop calculations of
Section~11.2.  It seems worthwhile to develop a formalism that will allow us
to use geometrical arguments about spontaneous symmetry breaking even at the
quantum level.

To define our goal somewhat better, consider the problem of determining the
vacuum expectation value of the quantum field $\phi$.  This expectation value
should be determined as a function of the parameters of the Lagrangian.  At
the classical level, it is easy to compute $\langle\phi\rangle$; one minimizes
the potential energy.  However, as we have seen in the previous section, this
classical value can be altered by perturbative loop corrections.  In fact, we
saw that $\langle\phi\rangle$ could be shifted by a potentially divergent
quantity, which we needed to control by renormalization.

It would be wonderful if, in the full quantum field theory, there were a
function whose minimum gave the exact value of $\langle\phi\rangle$.  This
function would agree with the classical potential energy to lowest order in
perturbation theory, but it would be modified in higher orders by quantum
corrections.  Presumably, these corrections would need renormalization to
remove infinities.  Nevertheless, after renormalization, this quantity should
give the same relations between $\langle\phi\rangle$ and particle masses and
couplings that we would find by direct Feynman diagram calculations.  In this
section, we will exhibit a function with these properties, called the
\emph{effective potential}.  In Section~11.4 we will explain how to compute
the effective potential in perturbation theory, in terms of renormalized
masses and couplings.  Then we will go on to use it as a tool in analyzing the
renormalizability of theories with hidden symmetry.

To identify the effective potential, consider the analogy between quantum
field theory and statistical mechanics set out in Section~9.3.  In that
section, we derived a correspondence between the correlation functions of a
quantum field and those of a related statistical system, with quantum
fluctuations being replaced by thermal fluctuations.  At zero temperature the
thermodynamic ground state is the state of lowest energy, but at nonzero
temperature we still have a geometrical picture of the preferred
thermodynamic state: It is the state that minimizes the Gibbs free energy.
More explicitly, taking the example of a magnetic system, one defines the
Helmholtz free energy $F(H)$ by
\begin{equation}
  Z(H) = e^{-\beta F(H)} = \int \mathcal{D}s\ \exp\Bigl[ -\beta
    \int d^3x\ \bigl( \mathcal{H}(s) - H s \bigr) \Bigr] .
  \label{eq:11.40}
\end{equation}
The thermodynamic state of the system at a given value of $H$ is the state
that minimizes the Gibbs free energy $G(M)$, which is related to the Helmholtz
free energy $F(H)$ by a Legendre transform,
\begin{equation}
  G(M) = F(H) - HM ,
\end{equation}
where $M = \langle s\rangle$ is the magnetization.  The magnetization is the
value of $s$ that minimizes $G(M)$.

In quantum field theory, we can define the analogues of these quantities.  The
generating functional is
\begin{equation}
  Z[J] = e^{-iE[J]} = \int \mathcal{D}\phi\ \exp\Bigl[ i\int d^4x\
    \bigl( \mathcal{L} + J\phi \bigr) \Bigr] ,
  \label{eq:11.43}
\end{equation}
where $E[J]$ is the analogue of the Helmholtz free energy.  The classical
field is defined by
\begin{equation}
  \phi_{cl}(x) = \langle\phi(x)\rangle_J = \frac{\delta E[J]}
    {\delta J(x)} .
  \label{eq:11.44}
\end{equation}
The \emph{effective action} $\Gamma[\phi_{cl}]$ is defined by the Legendre
transform
\begin{equation}
  \Gamma[\phi_{cl}] = -E[J] - \int d^4x\ J(x)\phi_{cl}(x) ,
  \label{eq:11.47}
\end{equation}
where $J$ is eliminated in favor of $\phi_{cl}$ through \eqref{eq:11.44}.
The functional derivative of $\Gamma$ is
\begin{equation}
  \frac{\delta\Gamma[\phi_{cl}]}{\delta\phi_{cl}(x)} = -J(x) .
  \label{eq:11.48}
\end{equation}
Thus, when $J = 0$, the effective action is extremized by the vacuum
expectation value of the field.  For translation-invariant field
configurations, we define the \emph{effective potential} by
\begin{equation}
  \Gamma[\phi_{cl}] = -\int d^4x\ V_{eff}(\phi_{cl}) ,
  \label{eq:11.50}
\end{equation}
for constant $\phi_{cl}$.  Then the vacuum expectation value of $\phi$ is the
value that minimizes $V_{eff}$:
\begin{equation}
  \frac{\partial V_{eff}}{\partial\phi_{cl}}\Big|_{\phi_{cl} =
    \langle\phi\rangle} = 0 .
  \label{eq:11.51}
\end{equation}
Each solution of Eq.~\eqref{eq:11.51} is a translation-invariant state with
$J = 0$.  Equation \eqref{eq:11.47} implies that $\Gamma = -E$ in this case,
and therefore that $V_{eff}(\phi_{cl})$, evaluated at a solution to
\eqref{eq:11.51}, is just the energy density of the corresponding state.

Figure~\ref{fig:11.6} illustrates one possible shape for the function
$V_{eff}(\phi)$.  The local maxima (or, for systems of several fields
$\phi^i$, possible saddle points) are unstable configurations that cannot be
realized as stationary states.  The figure also contains a local minimum of
$V_{eff}$ that is not the absolute minimum; this is a \emph{metastable vacuum
state}, which can decay to the true vacuum by quantum-mechanical tunneling.
The absolute minimum of $V_{eff}$ is the state of lowest energy in the
theory, and thus the true, stable, vacuum state.  A system with spontaneously
broken symmetry will have several minima of $V_{eff}$, all with the same
energy by virtue of the symmetry.  The choice of one among these vacua is the
spontaneous symmetry breaking.

\begin{figure}[t]
  \centering
  \includegraphics[width=0.55\textwidth]{fig11.6.png}
  \caption{A possible form for the effective potential in a scalar field
    theory.  The extrema of the effective potential occur at the points
    $\phi_{cl} = \phi_1, \phi_2, \phi_3$.  The true vacuum state is the one
    corresponding to $\phi_1$.  The state $\phi_2$ is unstable.  The state
    $\phi_3$ is metastable, but it can decay to $\phi_1$ by
    quantum-mechanical tunneling.}
  \label{fig:11.6}
\end{figure}

\begin{figure}[t]
  \centering
  \includegraphics[width=0.55\textwidth]{fig11.7.png}
  \caption{Exact convex form of the effective potential for the system of
    Fig.~\ref{fig:11.6}.}
  \label{fig:11.7}
\end{figure}

In drawing Fig.~\ref{fig:11.6}, we have assumed that we are computing the
effective potential for a fixed constant background value of $\phi$.  Under
some circumstances, this state does not give the true minimum energy
configuration for states with a given expectation value of $\phi$.  This
mismatch can occur in the following way: In a system for which the effective
potential for constant background fields is given by Fig.~\ref{fig:11.6},
consider choosing a value of $\phi_{cl}$ that is intermediate between the
locally stable vacuum states $\phi_1$ and $\phi_3$:
\begin{equation}
  \phi_{cl} = x\phi_1 + (1-x)\phi_3, \qquad 0 < x < 1 .
  \label{eq:11.52}
\end{equation}
The assumption of a constant background field gives a large value of the
effective potential, as indicated in the figure.  We can obtain a lower-energy
configuration by considering states with macroscopic regions in which
$\langle\phi\rangle = \phi_1$ and other regions in which
$\langle\phi\rangle = \phi_3$, in such a way that the average value of
$\langle\phi\rangle$ over the whole system is $\phi_{cl}$.  For such a
configuration, the average vacuum energy is given by
\begin{equation}
  V_{eff}(\phi_{cl}) = xV_{eff}(\phi_1) + (1-x)V_{eff}(\phi_3) ,
  \label{eq:11.53}
\end{equation}
as shown in Fig.~\ref{fig:11.7}.  The interpolation \eqref{eq:11.53} is the
field theoretic analogue of the Maxwell construction for the thermodynamic
free energy.  In general, for any $\phi_{cl}$, $\phi_1$, $\phi_3$ satisfying
\eqref{eq:11.52}, the estimate \eqref{eq:11.53} will be an upper bound to the
effective potential; we say that the effective potential is a \emph{convex
function} of $\phi_{cl}$.%
\footnote{The convexity of the Gibbs free energy is a well-known exact result
in statistical mechanics; see, for example, D.~Ruelle, \emph{Statistical
Mechanics} (W.~A.~Benjamin, Reading, Mass., 1969).}

Just as in thermodynamics, straightforward schemes for computing the
effective potential do not take account of the possibility of phase separation
and so lead to a structure of unstable and metastable configurations of the
type shown in Fig.~\ref{fig:11.6}.  The Maxwell construction must be performed
by hand to yield the final form of $V_{eff}(\phi_{cl})$.  Fortunately, the
absolute minimum of $V_{eff}$ is not affected by this nicety.

We have now solved the problem that we posed at the beginning of this
section: The effective potential, defined by Eqs.~\eqref{eq:11.47} and
\eqref{eq:11.50}, gives an easily visualized function whose minimization
defines the exact vacuum state of the quantum field theory, including all
effects of quantum corrections.  It is not obvious from these definitions how
to compute $V_{eff}$.  We will see how to do so in the next section, by direct
evaluation of the functional integral.

\section{Computation of the Effective Action}

Now that we have defined the object whose minimization gives the exact vacuum
state of a quantum field theory, we must learn how to compute it.  This can be
done in more than one way.  The simplest method, which we will use here,
requires that we be bold enough to evaluate the complete effective action
$\Gamma$ directly from its functional integral definition.%
\footnote{This method is due to R.~Jackiw, \emph{Phys.\ Rev.} \textbf{D9},
1686 (1974).}  After computing $\Gamma$, we can obtain $V_{eff}$ by
specializing to constant values of $\phi_{cl}$.

Our plan is to find a perturbation expansion for the generating functional
$Z$, starting with its functional integral definition \eqref{eq:11.43}.  We
will then take the logarithm to obtain the energy functional $E$, and finally
Legendre-transform according to Eq.~\eqref{eq:11.47} to obtain $\Gamma$.  We
will use renormalized perturbation theory, so it is convenient to split the
Lagrangian as we did in Eq.~\eqref{eq:10.18}, into a piece depending on
renormalized parameters and one containing the counterterms:
\begin{equation}
  \mathcal{L} = \mathcal{L}_1 + \delta\mathcal{L} .
  \label{eq:11.54}
\end{equation}
We wish to compute $\Gamma$ as a function of $\phi_{cl}$.  But the functional
$Z[J]$ depends on $\phi_{cl}$ through its dependence on $J$.  Thus, we must
find, at least implicitly, a relation between $J(x)$ and $\phi_{cl}(x)$.  At
the lowest order in perturbation theory, that relation is just the classical
field equation:
\begin{equation}
  0 = \phi_{cl} - \frac{\delta}{\delta J} \cdots
  \qquad\text{(to lowest order)} .
\end{equation}
Let us define $J_1(x)$ to be whatever function satisfies this equation
exactly, when $\mathcal{L} = \mathcal{L}_1$:
\begin{equation}
  0 = \frac{\delta\mathcal{L}_1}{\delta\phi}\Big|_{\phi=\phi_{cl}}
    + J_1(x) \qquad \text{(exactly)} .
  \label{eq:11.55}
\end{equation}
We will think of the difference between $J$ and $J_1$ as a counterterm,
analogous to $\delta\mathcal{L}$, so we write
\begin{equation}
  J(x) = J_1(x) + \delta J(x) ,
  \label{eq:11.56}
\end{equation}
where $\delta J$ is determined, order by order in perturbation theory, by the
original definition \eqref{eq:11.44} of $\phi_{cl}$, namely
$\langle\phi(x)\rangle_{J=0} = \phi_{cl}(x)$.

Using this notation, we rewrite Eq.~\eqref{eq:11.43} as
\begin{equation}
  Z[J] = \int \mathcal{D}\phi\ \exp\Bigl[ i\int d^4x\ \bigl(
    \mathcal{L}_1 + J_1\phi \bigr)\ \Bigr]\ \exp\Bigl[ i\int d^4x\ \bigl(
      \delta\mathcal{L} + \delta J\phi \bigr) \Bigr] .
  \label{eq:11.57}
\end{equation}
The second exponential contains the counterterms; leave this aside for the
moment.  In the first exponential, expand the exponent about $\phi_{cl}$ by
replacing $\phi(x) = \phi_{cl}(x) + \eta(x)$.  This exponent takes the form
\begin{equation}
  \begin{aligned}
    \int d^4x\ \bigl( \mathcal{L}_1 + J_1\phi \bigr)
    &= \int d^4x\ \bigl( \mathcal{L}_1[\phi_{cl}] + J_1\phi_{cl} \bigr)
      + \int d^4x\ \eta(x)\Bigl( \frac{\delta\mathcal{L}_1}{\delta\phi}
        + J_1 \Bigr)\\
    &\quad + \frac{1}{2}\int d^4x\,d^4y\ \eta(x)\eta(y)\
      \frac{\delta^2\mathcal{L}_1}{\delta\phi(x)\delta\phi(y)}
      + \cdots ,
  \end{aligned}
  \label{eq:11.58}
\end{equation}
where the various functional derivatives of $\mathcal{L}_1$ are evaluated at
$\phi_{cl}(x)$.  Notice that the term linear in $\eta$ vanishes by the use of
Eq.~\eqref{eq:11.55}.  The integral over $\eta$ is thus a Gaussian integral,
with the cubic and higher terms giving perturbative corrections.

We will describe a formal evaluation of this integral, following the
prescriptions of Section~9.2.  The ingredients in this evaluation will be the
coefficients of Eq.~\eqref{eq:11.58}, that is, the successive functional
derivatives of $\mathcal{L}_1$.  For the moment, please accept that these give
well-defined operators.  After presenting a general expression for
$\Gamma[\phi_{cl}]$, we will carry out this calculation explicitly in a scalar
field theory example.  We will see in this example that the formal operators
correspond to expressions familiar from Feynman diagram perturbation theory.

Let us, then, consider performing the integral over $\eta(x)$ using the
expansion \eqref{eq:11.58}.  Keeping only the terms up to quadratic order in
$\eta$, and still neglecting the counterterms, we have a pure Gaussian
integral, which can be evaluated in terms of a functional determinant:
\begin{equation}
  \int \mathcal{D}\eta\ \exp\Bigl[ \frac{i}{2}\int d^4x\,d^4y\
    \eta(x)\,\frac{\delta^2\mathcal{L}_1}{\delta\phi(x)\delta\phi(y)}\,
    \eta(y) \Bigr] = \Bigl[ \det\Bigl( -i\frac{\delta^2\mathcal{L}_1}
      {\delta\phi(x)\delta\phi(y)} \Bigr) \Bigr]^{-1/2} .
  \label{eq:11.59}
\end{equation}
This functional determinant will give us the lowest-order quantum correction
to the effective action, and for many purposes it is unnecessary to go further
in the expansion \eqref{eq:11.58}.  Later we will see that if we do include
the cubic and higher terms in $\eta$, these produce a Feynman diagram
expansion of the functional integral \eqref{eq:11.57} in which the propagator
is the operator inverse
\begin{equation}
  \Bigl( -i\frac{\delta^2\mathcal{L}_1}{\delta\phi(x)\delta\phi(y)}
    \Bigr)^{-1} ,
  \label{eq:11.60}
\end{equation}
and the vertices are the third and higher functional derivatives of
$\mathcal{L}_1$.

Finally, let us put back the effects of the second exponential in
Eq.~\eqref{eq:11.57}, that is, the counterterm Lagrangian.  It is useful to
expand this term about $\phi = \phi_{cl}$, writing it as
\begin{equation}
  \bigl( \delta\mathcal{L}[\phi_{cl}] + \delta J\phi_{cl} \bigr)
  + \bigl( \delta\mathcal{L}[\phi_{cl}+\eta] - \delta\mathcal{L}[\phi_{cl}]
    + \delta J\eta \bigr) .
  \label{eq:11.61}
\end{equation}
The second term of \eqref{eq:11.61} can be expanded as a Taylor series in
$\eta$; the successive terms give counterterm vertices which can be included
in the aforementioned Feynman diagrams.  The first term is a constant with
respect to the functional integral over $\eta$, and therefore gives additional
terms in the exponent of Eq.~\eqref{eq:11.59}.

Combining the integral \eqref{eq:11.59} with the contributions from
higher-order vertices and counterterms, one can obtain a complete expression
for the functional integral \eqref{eq:11.57}.  We will see in the example
below that the Feynman diagrams representing the higher-order terms can be
arranged to give the exponential of the sum of connected diagrams.  Thus one
obtains the following expression for $E[J]$:
\begin{equation}
  \begin{aligned}
    -iE[J] &= i\int d^4x\ \bigl( \mathcal{L}_1[\phi_{cl}] + J\phi_{cl} \bigr)
      - \frac{i}{2}\log\det\Bigl[ -i\frac{\delta^2\mathcal{L}_1}
        {\delta\phi(x)\delta\phi(y)} \Bigr]\\
    &\quad + (\text{connected diagrams}) + i\int d^4x\ \bigl(
      \delta\mathcal{L}[\phi_{cl}] + \delta J\phi_{cl} \bigr) .
  \end{aligned}
  \label{eq:11.62}
\end{equation}
From this equation, $\Gamma$ follows directly: Using $\delta E/\delta\phi_{cl}
+ \delta J = J$ and the Legendre transform \eqref{eq:11.47}, we find
\begin{equation}
  \Gamma[\phi_{cl}] = \int d^4x\ \mathcal{L}_1[\phi_{cl}]
    + \frac{i}{2}\log\det\Bigl[ -i\frac{\delta^2\mathcal{L}_1}
      {\delta\phi(x)\delta\phi(y)} \Bigr]
    - i\cdot(\text{connected diagrams}) + \int d^4x\ \delta\mathcal{L}[\phi_{cl}].
  \label{eq:11.63}
\end{equation}
Notice that there are no terms remaining that depend explicitly on $J$; thus,
$\Gamma$ is expressed as a function of $\phi_{cl}$, as it should be.  The
Feynman diagrams contributing to $\Gamma[\phi_{cl}]$ have no external lines,
and the simplest ones turn out to have two loops.  The lowest-order quantum
correction to $\Gamma$ is given by the functional determinant, and this term
is all that we will make use of in this book.

The last term of \eqref{eq:11.63} provides a set of counterterms that can be
used to satisfy the renormalization conditions on $\Gamma$ and, in the
process, to cancel divergences that appear in the evaluation of the functional
determinant and the diagrams.  We will show in the example below exactly how
this cancellation works.  The renormalization conditions will determine all of
the counterterms in $\delta\mathcal{L}$.  However, the formalism we have
constructed contains a new counterterm $\delta J$.  That coefficient is
determined by the following special criterion: In Eq.~\eqref{eq:11.55}, we
set up our analysis in such a way that, at the leading order,
$\langle\phi\rangle = \phi_{cl}$.  Potentially, however, this relation could
break down at higher orders: The quantity $\langle\phi\rangle$ could receive
additional contributions from Feynman diagrams that might shift it from the
value $\phi_{cl}$.  This will happen if there are nonzero tadpole diagrams
that contribute to $\langle\eta\rangle$.  But this amplitude also receives a
contribution from the counterterm $(\delta J\eta)$ in \eqref{eq:11.61}.  Thus
we can maintain $\langle\eta\rangle = 0$, and in the process determine
$\delta J$ to any order, by adjusting $\delta J$ to satisfy the diagrammatic
equation
\begin{equation}
  \delta J\ \bigl( \text{tadpole diagram} \bigr)
  = -\bigl( \text{sum of all tadpole diagrams} \bigr) = 0 .
  \label{eq:11.64}
\end{equation}
In practice, we will satisfy this condition by simply ignoring any
one-particle-irreducible one-point diagram, since any such diagram will be
canceled by adjustment of $\delta J$.  The removal of these tadpole diagrams,
which we needed some effort to arrange in Section~11.2, is thus built in here
as a natural part of the formalism.

\subsection{The Effective Action in the Linear Sigma Model}

In Eq.~\eqref{eq:11.63}, we have given a complete, though not exactly
transparent, evaluation of $\Gamma[\phi_{cl}]$.  Let us now clarify the
meaning of this equation, and also put it to some good use, by computing
$\Gamma[\phi_{cl}]$ in the linear sigma model.  We will see that the results
that we obtained by brute-force perturbation theory in Section~11.2 emerge
much more naturally from Eq.~\eqref{eq:11.63}.

We begin again with the Lagrangian \eqref{eq:11.5}:
\begin{equation}
  \mathcal{L} = \tfrac{1}{2}(\partial_\mu\phi^i)^2 + \tfrac{1}{2}\mu^2
    (\phi^i)^2 - \frac{\lambda}{4}\bigl[ (\phi^i)^2 \bigr]^2 .
  \label{eq:11.65}
\end{equation}
Expand about the classical field: $\phi^i = \phi^i_{cl} + \eta^i$.  Because we
expect to find a translation-invariant vacuum state, we will specialize to the
case of a constant classical field.  This will simplify some elements of the
calculation below.  In particular, according to Eq.~\eqref{eq:11.50}, the
final result will be proportional to the four-dimensional volume $(VT)$ of
the functional integration.  When this dependence is factored out, we will
obtain a well-defined intensive expression for the effective potential.  In
any event, after this simplification, \eqref{eq:11.65} takes the form
\begin{equation}
  \mathcal{L} = \frac{1}{2}\bigl( \mu^2(\phi^i_{cl})^2
    - \frac{\lambda}{2}(\phi^i_{cl})^2)^2 \bigr)
    + \frac{1}{2}\eta^i\bigl( -\partial^2\delta^{ij} + (\lambda(\phi_{cl})^2
      - \mu^2)\delta^{ij} + 2\lambda\phi^i_{cl}\phi^j_{cl} \bigr)\eta^j
    + \cdots .
  \label{eq:11.66}
\end{equation}
The quadratic term defines the operator whose determinant we must evaluate.
For the case of a constant classical field, we can choose the coordinate
system so that $\phi_{cl}$ points along the $N$th direction, as in
\eqref{eq:11.7}: $\phi^i_{cl} = (0, \ldots, 0, \phi_{cl})$.  Then the operator
is diagonal, with eigenvalues
\begin{equation}
  m^2_\sigma = \mu^2 - 3\lambda\phi_{cl}^2
  \qquad\text{and}\qquad
  m^2_\pi = \mu^2 - \lambda\phi_{cl}^2 ,
  \label{eq:11.69}
\end{equation}
for the $\sigma$ and $\pi$ fields, respectively.  Note that the $\pi$
eigenvalue vanishes when $\phi_{cl}^2 = \mu^2/\lambda$, the classical vacuum
expectation value.  This is Goldstone's theorem at the level of the effective
potential: the mass matrix has a zero eigenvalue precisely at the minimum.

The functional determinant in \eqref{eq:11.63} is then
\begin{equation}
  \frac{i}{2}\log\det\Bigl[ -i\frac{\delta^2\mathcal{L}_1}
    {\delta\phi(x)\delta\phi(y)} \Bigr]
  = \frac{i}{2}\log\det\bigl[ -i(\partial^2 + M^2) \bigr],
  \label{eq:11.71}
\end{equation}
where $M^2$ is the mass matrix with eigenvalues given by \eqref{eq:11.69}.
Using the identity $\log\det M = \mathrm{tr}\log M$ and evaluating the
resulting Gaussian functional integral, we obtain
\begin{equation}
  V^{(1)}_{eff}(\phi_{cl}) = \frac{1}{2(4\pi)^{d/2}}\ \Gamma(-d/2)\
    \Bigl[ (m_\sigma^2)^{d/2} + (N-1)(m_\pi^2)^{d/2} \Bigr] .
  \label{eq:11.73}
\end{equation}
Adding the counterterm Lagrangian as in Eq.~\eqref{eq:11.15}, we find
\begin{equation}
  \begin{aligned}
    V_{eff}(\phi_{cl}) &= -\frac{1}{2}\mu^2\phi_{cl}^2 + \frac{\lambda}{4}
      \phi_{cl}^4 + \frac{1}{2(4\pi)^{d/2}}\ \Gamma(-d/2)\ \Bigl[
        \bigl( \mu^2 - 3\lambda\phi_{cl}^2 \bigr)^{d/2}\\
    &\quad + (N-1)\bigl( \mu^2 - \lambda\phi_{cl}^2 \bigr)^{d/2} \Bigr]
      + 2\delta_\mu\phi_{cl}^2 + 4\delta_\lambda\phi_{cl}^4 .
  \end{aligned}
  \label{eq:11.74}
\end{equation}
Here we have written $(\phi_{cl})^2$ as a shorthand for
$(\phi^i_{cl})^2$.  Since the second line of this result is the leading
radiative correction, we might expect that the result has the structure of a
one-loop Feynman diagram.  Indeed, we see that this expression contains Gamma
functions and ultraviolet divergences similar to those that we found in the
one-loop computations of Section~11.2.  We will show below that this term in
fact has exactly the same ultraviolet divergences that we found in
Section~11.2.  These divergences will be subtracted by the counterterms in
the last line of Eq.~\eqref{eq:11.74}.

\begin{figure}[t]
  \centering
  \includegraphics[width=0.5\textwidth]{fig11.8.png}
  \caption{Feynman diagrams contributing to the evaluation of the effective
    potential of the $O(N)$ linear sigma model: (a) a diagram that is removed
    by \eqref{eq:11.64}; (b) the first nonzero diagrammatic corrections.}
  \label{fig:11.8}
\end{figure}

Since the computation of the determinant in Eq.~\eqref{eq:11.63} gives the
effect of one-loop corrections, we might expect the Feynman diagrams that
contribute to Eq.~\eqref{eq:11.63} to begin in two-loop order.  We can see
this explicitly for the case of the $O(N)$ sigma model.  The perturbation
expansion described below Eq.~\eqref{eq:11.60} involves the propagator that
is the inverse of Eq.~(11.67):
\begin{equation}
  \langle \eta^i(k)\eta^j(-k) \rangle = \frac{i\delta^{ij}}{k^2 - m^2} ,
  \label{eq:11.75}
\end{equation}
where $m^2$ is given by \eqref{eq:11.69}.  The vertices are given by the
terms of order $\eta^3$ and $\eta^4$ in the expansion of the Lagrangian.
Combining these ingredients, we find that the leading Feynman diagrams
contributing to the vacuum energy have the forms shown in Fig.~\ref{fig:11.8}.
The diagram of Fig.~\ref{fig:11.8}(a) is actually canceled by the effects of
the counterterm $\delta J$, as shown in Eq.~\eqref{eq:11.64}.  Thus the
leading diagrammatic contribution to the effective potential comes from the
two-loop diagrams of Fig.~\ref{fig:11.8}(b).

The result \eqref{eq:11.74} is manifestly $O(N)$-symmetric.  From the
question that we posed at the beginning of Section~11.2, we might have feared
that this property would be destroyed when we compute radiative corrections
about a state with spontaneously broken symmetry.  But $V_{eff}(\phi_{cl})$
is the function that we minimize to find the vacuum state, and so it should
properly be symmetric, even if the lowest-energy vacuum is asymmetric.  In
the formalism we have constructed here, there is no need to worry.  Formula
\eqref{eq:11.63} is manifestly invariant, term by term, under the original
$O(N)$ symmetry of the Lagrangian.  Thus we must necessarily have arrived at
an $O(N)$-symmetric result for $\Gamma_{eff}(\phi_{cl})$.

Before going on to determine $\delta_\mu$ and $\delta_\lambda$ precisely, we
might first check that the counterterms in Eq.~\eqref{eq:11.74} are sufficient
to make the expression for $\Gamma[\phi_{cl}]$ finite.  The factor
$\Gamma(-d/2)$ has poles at $d = 0, 2, 4$.  The pole at $d = 0$ is a constant,
independent of $\phi_{cl}$, and therefore without physical significance.  The
pole at $d = 2$ is an even quadratic polynomial in $\phi_{cl}$.  The pole at
$d = 4$ is an even quartic polynomial in $\phi_{cl}$.  Thus Eq.~\eqref{eq:11.74}
becomes a finite expression in the limit $d\to 2$ if we set
\begin{equation}
  \delta_\mu = -\lambda\mu^2(N+2)\ \frac{\Gamma(-d/2)}{(4\pi)^{d/2}}
    + \text{finite} .
\end{equation}
The expression is finite as $d\to 4$ if we set
\begin{equation}
  \delta_\mu = -\lambda\mu^2(N+2)\ \frac{\Gamma(2-d/2)}{(4\pi)^2}
    + \text{finite};\qquad
  \delta_\lambda = \lambda^2(N+8)\ \frac{\Gamma(2-d/2)}{(4\pi)^2}
    + \text{finite} .
  \label{eq:11.76}
\end{equation}
These expressions agree with our earlier results from Section~11.2,
Eqs.~\eqref{eq:11.33} and \eqref{eq:11.22}, in the limits $d\to 2$ and
$d\to 4$ respectively.

The finite parts of $\delta_\lambda$ and $\delta_\mu$ depend on the exact form
of the renormalization conditions that are imposed.  For example, in
Section~11.2, we imposed the condition \eqref{eq:11.16} that the vacuum
expectation value of $\phi$ equals $\mu/\sqrt{\lambda}$, and the additional
conditions in \eqref{eq:11.17} on the scattering amplitude and field strength
of the $\sigma$.  Condition \eqref{eq:11.16} is readily expressed in terms of
the effective potential as
\begin{equation}
  \frac{dV_{eff}}{d\phi_{cl}}\Big|_{\phi_{cl} = \mu/\sqrt{\lambda}} = 0 .
\end{equation}
Using the connection between derivatives of $\Gamma$ and
one-particle-irreducible amplitudes, we could write the other two conditions
as Fourier transforms to momentum space of functional derivatives of
$\Gamma[\phi_{cl}]$.  In this way, it is possible in principle to reconstruct
the particular renormalization scheme used in Section~11.2.

However, if we want to visualize the modification of the lowest-order results
that is induced by the quantum corrections, we can apply a renormalization
scheme that can be implemented more easily.  One such scheme, known as
\emph{minimal subtraction} (MS), is simply to remove the $(1/\epsilon)$ poles
(for $\epsilon = 4-d$) in potentially divergent quantities.  Normally, though,
these $(1/\epsilon)$ poles are accompanied by terms involving $\gamma$ and
$\log(4\pi)$.  It is convenient, and no more arbitrary, to subtract these
terms as well.  In this prescription, known as \emph{modified minimal
subtraction} or $\overline{MS}$ (``em-ess-bar''), one replaces
\begin{equation}
  \frac{1}{(4\pi)^{d/2}(m^2)^{2-d/2}}\ \Gamma(2-d/2)
  \to \frac{1}{(4\pi)^2}\ \Bigl( \frac{2}{\epsilon} - \gamma_E
    + \log 4\pi - \log m^2 \Bigr)
  \to \frac{1}{(4\pi)^2}\ \Bigl( -\log\Bigl( \frac{m^2}{M^2} \Bigr)
    + \cdots \Bigr),
  \label{eq:11.77}
\end{equation}
where $M$ is an arbitrary mass parameter that we have introduced to make the
final equation dimensionally correct.  You should think of $M$ as
parametrizing a sequence of possible renormalization conditions.  The
$\overline{MS}$ renormalization scheme usually puts one-loop corrections in an
especially simple form.  The price of this simplicity is that it normally
takes some effort to express physically measurable quantities in terms of the
parameters of the $\overline{MS}$ expression.

To apply the $\overline{MS}$ renormalization prescription to \eqref{eq:11.74},
we need to expand the divergent terms in this equation in powers of
$\epsilon$.  As an example, consider the $\overline{MS}$ regularization of
expression \eqref{eq:11.73}:
\begin{equation}
  \begin{aligned}
    \frac{1}{2(4\pi)^{d/2}}\ \Gamma(-d/2)\,m^{d/2}
    &= \frac{1}{(4\pi)^{d/2}}\ \frac{\Gamma(2-d/2)}{(d/2)(d/2-1)}\ m^2\
      (m^2)^{d/2-2}\\
    &= \frac{m^2}{2(4\pi)^2}\ \Bigl( -\frac{2}{\epsilon} - \log\Bigl(
      \frac{m^2}{M^2} \Bigr) + \frac{3}{2} \Bigr) .
  \end{aligned}
  \label{eq:11.78}
\end{equation}
Modifying our result \eqref{eq:11.74} in this way, we find
\begin{equation}
  \begin{aligned}
    V_{eff} &= -\frac{1}{2}\mu^2\phi_{cl}^2 + \frac{\lambda}{4}\phi_{cl}^4
      + \frac{1}{64\pi^2}\ \Bigl\{ (N-1)\bigl( \mu^2 - \lambda\phi_{cl}^2
        \bigr)^2\Bigl[ \log\Bigl( \frac{\mu^2 - \lambda\phi_{cl}^2}{M^2}
          \Bigr) - \frac{3}{2} \Bigr]\\
    &\quad + \bigl( \mu^2 - 3\lambda\phi_{cl}^2 \bigr)^2\Bigl[ \log\Bigl(
      \frac{\mu^2 - 3\lambda\phi_{cl}^2}{M^2} \Bigr) - \frac{3}{2}
      \Bigr] \Bigr\} .
  \end{aligned}
  \label{eq:11.79}
\end{equation}
The effective potential is thus modified to be slightly steeper at large
values of $\phi_{cl}$ and more negative at smaller values, as shown in
Fig.~\ref{fig:11.9}.  For each set of values of $\mu$, $\lambda$, and $M$, we
can determine the preferred vacuum state by minimizing $V_{eff}(\phi)$ with
respect to $\phi_{cl}$.  The correction to $V_{eff}$ is undefined when the
arguments of the logarithms become negative, but fortunately the minima of
$V_{eff}$ occur outside of this region, as is illustrated in the figure.

\begin{figure}[t]
  \centering
  \includegraphics[width=0.5\textwidth]{fig11.9.png}
  \caption{The effective potential for $\phi^4$ theory ($N=1$), with quantum
    corrections included as in Eq.~\eqref{eq:11.79}.  The lighter-weight
    curve shows the classical potential energy, for comparison.}
  \label{fig:11.9}
\end{figure}

Before going on, we would like to raise two questions about this expression
for the effective potential.  The problems that we will raise occur generically
in quantum field theory calculations, but expression \eqref{eq:11.79} provides
a concrete illustration of these difficulties.  Most of our discussion in the
next two chapters will be devoted to building a formalism within which these
questions can be answered.

First, it is troubling that, while our classical Lagrangian contained only
two parameters, $\mu$ and $\lambda$, the result \eqref{eq:11.79} depends on
three parameters, of which one is the arbitrary mass scale $M$.  A superficial
reply to this complaint can be given as follows: Consider the change in
$V_{eff}(\phi_{cl})$ that results from changing the value of $M^2$ to
$M^2 + \delta M^2$.  From the explicit form of \eqref{eq:11.79}, we can see
that this change is compensated completely by shifting the values of $\mu$ and
$\lambda$, according to
\begin{equation}
  \mu^2 \to \mu^2\Bigl( 1 + \frac{\lambda(N+2)}{16\pi^2}\ \log\Bigl(
    \frac{M^2}{M_0^2} \Bigr) \Bigr),\qquad
  \lambda \to \lambda\Bigl( 1 + \frac{\lambda(N+8)}{16\pi^2}\ \log\Bigl(
    \frac{M^2}{M_0^2} \Bigr) \Bigr) .
  \label{eq:11.80}
\end{equation}
Thus, a change in $M^2$ is completely equivalent to changes in the parameters
$\mu$ and $\lambda$.  It is not clear, however, why this should be true or how
this fact helps us understand the dependence of our formulae on $M^2$.

The second problem arises from the fact that the one-loop correction in
Eq.~\eqref{eq:11.79} includes a logarithm that can become large enough to
compensate the small coupling constant $\lambda$.  The problem is particularly
clear in the limit $\mu^2\to 0$; then Eq.~\eqref{eq:11.79} takes the form
\begin{equation}
  \begin{aligned}
    V_{eff} &= \frac{\lambda}{4}\phi_{cl}^4 + \frac{1}{64\pi^2}\ \Bigl[
      \lambda^2\phi_{cl}^4\Bigl( (N+8)\Bigl( \log\Bigl(
        \frac{\lambda\phi_{cl}^2}{M^2} \Bigr) - 1 \Bigr) + 9\log 3 \Bigr)
      \Bigr]\\
    &= \frac{\lambda}{4}\phi_{cl}^4\ \Bigl[ 1 + \frac{\lambda}{16\pi^2}\ \Bigl(
      (N+8)\Bigl( \log\Bigl( \frac{\lambda\phi_{cl}^2}{M^2} \Bigr)
        - 1 \Bigr) + 9\log 3 \Bigr) \Bigr] .
  \end{aligned}
  \label{eq:11.81}
\end{equation}
Where is the minimum of this potential?  If we take this expression at face
value, we find that $V_{eff}(\phi_{cl})$ passes through zero when $\phi_{cl}$
reaches a very small value of order
\begin{equation}
  \phi_{cl} \sim \frac{M}{\sqrt{\lambda}}\ \exp\Bigl[ -\frac{16\pi^2}
    {(N+8)\lambda} \Bigr] ,
\end{equation}
and, near this point, attains a minimum with a nonzero value of $\phi_{cl}$.
But the zero occurs by the cancellation of the leading term against the
quantum correction.  In other words, perturbation theory breaks down
completely before we can address the question of whether $V_{eff}(\phi_{cl})$,
for $\mu^2 = 0$, has a symmetry-breaking minimum.  It seems that our present
tools are quite inadequate to resolve this case.

Although it is far from obvious, these two problems turn out to be related to
each other.  One of our major results in Chapter~12 will be an explanation of
the interrelation of $M^2$, $\lambda$, and $\mu^2$ displayed in
Eq.~\eqref{eq:11.80}.  Then, in Chapter~13, we will use the insight we have
gained from this analysis to solve completely the second problem of the
appearance of large logarithms.  Before beginning that study, however, there
are a few issues we have yet to discuss in the more formal aspects of the
renormalization of theories with spontaneously broken symmetry.

\section{The Effective Action as a Generating Functional}

Now that we have defined the effective action and computed it for one
particular theory, let us return to our goal of understanding the
renormalization of theories with hidden symmetry.  In Section~11.6 we will use
the effective action as a tool in achieving this goal.  First, however, we
must investigate in more detail the relation between the effective action and
Feynman diagrams.

We saw in Section~9.2 that the functional derivatives of $Z[J]$ with respect
to $J(x)$ produce the correlation functions of the scalar field (see, for
example, Eq.~\eqref{eq:9.35}).  In other words, $Z[J]$ is the generating
functional of correlation functions.  Our goal now is to show that
$\Gamma[\phi_{cl}]$ is also such a generating functional; specifically, it is
the generating functional of \emph{one-particle-irreducible} (1PI)
correlation functions.  Since the 1PI correlation functions figure
prominently in the theory of renormalization, this result will be central in
the discussion of renormalization in the following section.

To begin, let us consider the functional derivatives not of $\Gamma[\phi_{cl}]$,
but of $E[J] = -i\log Z[J]$.  The first derivative, given in Eq.
\eqref{eq:11.44}, is precisely $\langle\phi(x)\rangle$.  The second
derivative is
\begin{equation}
  \frac{\delta^2 E[J]}{\delta J(x)\delta J(y)}
  = -i\Bigl( \langle\phi(x)\phi(y)\rangle
    - \langle\phi(x)\rangle\langle\phi(y)\rangle \Bigr) ,
  \label{eq:11.82}
\end{equation}
which is the connected two-point function (up to the factor $-i$).  The
higher derivatives of $E[J]$ give the connected $n$-point functions:
\begin{equation}
  \frac{\delta^n E[J]}{\delta J(x_1)\cdots\delta J(x_n)} = -i\,
    \langle\phi(x_1)\cdots\phi(x_n)\rangle_{\mathrm{connected}} .
  \label{eq:11.83}
\end{equation}
Thus, $E[J]$ is the generating functional of connected correlation functions.

Now, the effective action is related to $E[J]$ by the Legendre transform
\eqref{eq:11.47}.  The first functional derivative of $\Gamma$ is
\begin{equation}
  \frac{\delta\Gamma[\phi_{cl}]}{\delta\phi_{cl}(x)} = -J(x) .
  \label{eq:11.84}
\end{equation}
The second derivative is related to the inverse of the connected two-point
function.  Differentiating \eqref{eq:11.44} with respect to $\phi_{cl}$, we
obtain
\begin{equation}
  \delta(x-y) = \int d^4z\ \frac{\delta^2 E[J]}{\delta J(x)\delta J(z)}\
    \frac{\delta J(z)}{\delta\phi_{cl}(y)}
  = \int d^4z\ \frac{\delta^2 E[J]}{\delta J(x)\delta J(z)}\
    \Bigl( -\frac{\delta^2\Gamma}{\delta\phi_{cl}(z)\delta\phi_{cl}(y)}
      \Bigr) .
  \label{eq:11.90}
\end{equation}
Using \eqref{eq:11.82}, this says that the second derivative of $\Gamma$ is
the inverse of the connected two-point function (with appropriate factors of
$i$).  Diagrammatically, we can write
\begin{equation}
  \frac{\delta^2\Gamma[\phi_{cl}]}{\delta\phi_{cl}(x)\delta\phi_{cl}(y)}
  = iD^{-1}(x-y) ,
  \label{eq:11.91}
\end{equation}
where $D(x-y)$ is the full propagator (the connected two-point function).
Thus, the second derivative of the effective action is the \emph{inverse
propagator}, that is, the full two-point function with the external
propagators removed.

For the three-point function, we can differentiate \eqref{eq:11.91} and use
the chain rule to obtain
\begin{equation}
  \frac{\delta^3 E[J]}{\delta J(x)\delta J(y)\delta J(z)}
  = i\int d^4u\,d^4v\,d^4w\ D(x-u)D(y-v)D(z-w)\
    \frac{\delta^3\Gamma[\phi_{cl}]}
      {\delta\phi_{cl}(u)\delta\phi_{cl}(v)\delta\phi_{cl}(w)} .
  \label{eq:11.95}
\end{equation}
This relation is more clearly expressed diagrammatically.  The left-hand side
is the connected three-point function.  If we extract exact propagators as
indicated in \eqref{eq:11.95}, this decomposes as follows: the third
derivative of $i\Gamma[\phi_{cl}]$ is just the connected correlation function
with all three full propagators removed, that is, the
one-particle-irreducible three-point function:
\begin{equation}
  \frac{\delta^3\Gamma[\phi_{cl}]}{\delta\phi_{cl}(x)\delta\phi_{cl}(y)
    \delta\phi_{cl}(z)} = \text{(1PI three-point function)} .
\end{equation}
By similar, if increasingly complicated, manipulations, one can derive the
same relation for each successive derivative of $\Gamma$.  For example,
differentiating Eq.~\eqref{eq:11.95}, we eventually find (using matrix
notation with repeated indices implicitly integrated over) that the fourth
derivative of $i\Gamma$ is the one-particle-irreducible four-point function.
The general relation (for $n \ge 3$) is
\begin{equation}
  \frac{\delta^n\Gamma[\phi_{cl}]}{\delta\phi_{cl}(x_1)\cdots
    \delta\phi_{cl}(x_n)} = \text{(1PI } n\text{-point function)} .
  \label{eq:11.96}
\end{equation}
In other words, the effective action is the generating functional of
one-particle-irreducible correlation functions.

This conclusion implies that $\Gamma$ contains the complete set of physical
predictions of the quantum field theory.  Let us review how this information
is encoded.  The vacuum state of the field theory is identified as the minimum
of the effective potential.  The location of the minimum determines whether
the symmetries of the Lagrangian are preserved or spontaneously broken.  The
second derivative of $\Gamma$ is the inverse propagator.  The poles of the
propagator, or the zeros of the inverse propagator, give the values of the
particle masses.  Thus the particle masses $m^2$ are determined as the values
of $p^2$ that solve the equation
\begin{equation}
  \int d^4x\ e^{ip\cdot(x-y)}\ \frac{\delta^2\Gamma}{\delta\phi(x)
    \delta\phi(y)} = 0 .
  \label{eq:11.97}
\end{equation}
The higher derivatives of $\Gamma$ are the one-particle-irreducible
amplitudes.  These can be connected by full propagators and joined together to
construct four- and higher-point connected amplitudes, which give the
$S$-matrix elements.  Thus, from the knowledge of $\Gamma$, we can reconstruct
the qualitative behavior of the quantum field theory, its pattern of
symmetry-breaking, and then the quantitative details of its particles and
their interactions.

\section{Renormalization and Symmetry: General Analysis}

In our analysis of the divergences of quantum field theories (especially in
the paragraph below Eq.~\eqref{eq:10.4}), we noted that the basic divergences
of Feynman integrals are associated with one-particle-irreducible diagrams.
Thus we might expect that the effective action will be a useful object in
discussing the renormalizability of quantum field theories, especially those
with spontaneously broken symmetry.  In this section we will make use of the
effective action in precisely this way.

In Section~11.4, we saw in a particular example that the formalism for
calculating the effective action provides the counterterms needed to remove
the ultraviolet divergences, at least at the one-loop level.  These
counterterms were exactly those of the original Lagrangian.  We will now argue
that this set of counterterms is always sufficient---to all orders and for any
renormalizable field theory---by applying the power-counting arguments of
Section~10.1 directly to the computation of the effective action.  We will use
the language of scalar field theories, but the arguments can be generalized to
theories of spinor and vector fields.

Consider first the computation of the effective potential for constant
($x$-independent) classical fields, in a field theory with an arbitrary number
of fields $\phi^i$.  The effective potential has mass dimension 4, so we expect
that $V_{eff}(\phi_{cl})$ will have divergent terms up to $\Lambda^4$.  To
understand these divergences, expand $V_{eff}(\phi_{cl})$ in a Taylor series:
\begin{equation}
  V_{eff}(\phi_{cl}) = A_0 + A_2\phi_{cl}^i\phi_{cl}^i
    + \frac{1}{4!}A_4(\phi_{cl}^i\phi_{cl}^i)^2 + \cdots .
\end{equation}
(In theories without a symmetry $\phi^i\to-\phi^i$, there might also be terms
linear and cubic in $\phi_{cl}$; we omit these for simplicity.)  The
coefficients $A_0$, $A_2$, $A_4$ have mass dimension, respectively, 4, 2, and
0; thus we expect them to contain $\Lambda^4$, $\Lambda^2$, and $\log\Lambda$
divergences, respectively.  The power-counting analysis predicts that all
higher terms in the Taylor series expansion should be finite.

The constant term $A_0$ is independent of $\phi_{cl}$; it has no physical
significance.  However, the divergences in $A_2$ and $A_4$ appear in physical
quantities, since these coefficients enter the inverse propagator
\eqref{eq:11.90} and the irreducible four-point function \eqref{eq:11.96} and
therefore appear in the computation of $S$-matrix elements.  There is one
further coefficient in the effective action that has nonnegative mass
dimension by power counting; this is the coefficient of the term quadratic in
$\partial_\mu\phi_{cl}$, which appears when the effective action is evaluated
for a nonconstant background field:
\begin{equation}
  \Delta\Gamma[\phi_{cl}] = \int d^4x\ B\ \frac{1}{2}(\partial_\mu\phi_{cl})^2 .
  \label{eq:11.98}
\end{equation}
All other coefficients in the Taylor expansion of the effective action in
powers of $\phi_{cl}$ are finite by power counting.

We can now argue that the counterterms of the original Lagrangian suffice to
remove the divergences that might appear in the computation of
$\Gamma[\phi_{cl}]$.  The argument proceeds in two steps.  We first use the
BPHZ theorem to argue that the divergences of Green's functions can be removed
by adjusting a set of counterterms corresponding to the possible operators
that can be added to the Lagrangian with coefficients of mass dimension
greater than or equal to zero.  The coefficients of these counterterms are in
1-to-1 correspondence with the coefficients $A_2$, $A_4$, and $B_2$ of the
effective action.  Next, we use the fact that the effective action is
manifestly invariant to the original symmetry group of the model.  This is
true even if the vacuum state of the model has spontaneous symmetry breaking.
This symmetry of the effective action follows from the analysis of
Section~11.4, since the method we presented there for computing the effective
action is manifestly invariant to the original symmetry of the Lagrangian.

Combining these two results, we conclude that the effective action can always
be made finite by adjusting the set of counterterms that are invariant to the
original symmetry of the theory, even if this symmetry is spontaneously
broken.  By using the results of Section~11.5, which explain how to construct
the Green's functions of the theory from the functional derivatives of the
effective action, this conclusion of renormalizability extends to all the
Green's functions of the theory.

To make this abstract argument more concrete, we will demonstrate in a simple
example how the functional derivatives of the effective action yield a set of
Feynman diagrams whose divergences correspond to symmetric counterterms.  Let
us, then, return once again to the $O(N)$-invariant linear sigma model and
compute the second functional derivative of $\Gamma[\phi_{cl}]$.  If the whole
formalism we have constructed hangs together, we should be able to recognize
the result as the Feynman diagram expansion of the inverse propagator, with
divergences corresponding to the counterterms of $O(N)$-symmetric scalar
field theory.

To begin, we write out expression \eqref{eq:11.63} explicitly for this model:
\begin{equation}
  \Gamma[\phi_{cl}] = \int d^4x\ \Bigl( \tfrac{1}{2}(\partial_\mu\phi_{cl})^2
    + \tfrac{1}{2}m^2(\phi_{cl})^2 - \frac{\lambda}{4}\bigl[ (\phi_{cl})^2
      \bigr]^2 \Bigr) + \frac{i}{2}\log\det[-iD^{ij}] + \cdots ,
  \label{eq:11.99}
\end{equation}
where
\begin{equation}
  -iD^{ij} = \partial^2\delta^{ij} + \bigl( \lambda(\phi_{cl}(x))^2 - \mu^2
    \bigr)\delta^{ij} + 2\lambda\phi^i_{cl}(x)\phi^j_{cl}(x) .
  \label{eq:11.100}
\end{equation}
For constant $\phi_{cl}$, $D^{ij}$ is the operator that, acting on a given
component of the scalar field, equals the Klein--Gordon operator with mass
squared given by Eq.~\eqref{eq:11.69}.  This is the leading-order
approximation to the inverse propagator of the linear sigma model.

To find the higher-order corrections to the inverse propagator, we must
compute the second functional derivative of the quantum correction terms in
$\Gamma[\phi_{cl}]$.  From \eqref{eq:11.99}, we find
\begin{equation}
  \begin{aligned}
    \frac{\delta^2\Gamma}{\delta\phi^i_{cl}(x)\delta\phi^j_{cl}(y)}
    &= iD^{ij}\delta(x-y) + \frac{\delta^2}{\delta\phi^i_{cl}(x)
      \delta\phi^j_{cl}(y)}\ \frac{i}{2}\log\det[-iD^{ij}] + \cdots .
  \end{aligned}
\end{equation}
The first term is just the Klein--Gordon operator $iD^{ij}\delta(x-y)$.  To
compute the second term, use identity \eqref{eq:9.78} for determinants of
matrices:
\begin{equation}
  \frac{\partial}{\partial a}\log\det M(a) = \frac{\partial}{\partial a}
    \mathrm{tr}\log M(a) = \mathrm{tr}\Bigl( M^{-1}\frac{\partial M}
      {\partial a} \Bigr) .
  \label{eq:11.101}
\end{equation}
Using this identity, we find
\begin{equation}
  \frac{\delta}{\delta\phi^k_{cl}(z)}\ \frac{i}{2}\log\det[-iD]
  = i\lambda\bigl( \phi^k_{cl}(z)\delta^{ij} + \phi^i_{cl}(z)\delta^{jk}
    + \phi^j_{cl}(z)\delta^{ik} \bigr)\ (D^{-1})^{ij}(z,z) .
  \label{eq:11.102}
\end{equation}
The quantity $(D^{-1})^{ij}(x,y)$ is the Klein--Gordon propagator.  To
differentiate a second time, we can use the identity Eq.~(11.94); this
yields
\begin{equation}
  \begin{aligned}
    &\frac{\delta^2}{\delta\phi^k_{cl}(z)\delta\phi^\ell_{cl}(w)}\
    \frac{i}{2}\log\det[-iD]\\
    &= -\lambda\bigl( \delta^{k\ell}\delta^{ij} + \delta^{ik}\delta^{j\ell}
      + \delta^{i\ell}\delta^{jk} \bigr)\ (D^{-1})^{ij}(z,z)\delta(z-w)\\
    &\quad + 2i\lambda^2\bigl( \phi^k_{cl}(z)\delta^{im} + \phi^i_{cl}(z)
      \delta^{mk} + \phi^m_{cl}(z)\delta^{ik} \bigr)\ (D^{-1})^{im}(z,w)\
      \bigl( \phi^\ell_{cl}(w)\delta^{nj} + \phi^n_{cl}(w)\delta^{n\ell}
        + \phi^{n}_{cl}(w)\delta^{\ell n} \bigr)\ (D^{-1})^{nj}(w,z) .
  \end{aligned}
  \label{eq:11.103}
\end{equation}
This is expected to be the formal correction to the inverse propagator at
one-loop order, and indeed we can recognize in \eqref{eq:11.103} the values of
the one-loop diagrams.  Notice how, in this derivation, every functional
derivative on $D^{-1}$ adds another propagator to the diagram and thus lowers
the degree of divergence, in conformity with our general arguments in
Section~10.1.

This example illustrates that the successive functional derivatives of
$\Gamma[\phi_{cl}]$ are computed by a Feynman diagram expansion, with
propagators and vertices that depend on the classical field.  When the
classical field is a constant, the propagators reduce to ordinary
Klein--Gordon propagators and so the BPHZ theorem applies.  All ultraviolet
divergences can be removed from all of the amplitudes obtained by
differentiating $\Gamma[\phi_{cl}]$ by the use of the most general set of
mass, vertex, and field-strength renormalizations.  At the same time, the
perturbation theory is manifestly invariant to the symmetry of the original
Lagrangian, and so the only divergences that appear---and thus the only
counterterms required---are those that respect this symmetry.  In general,
then, all amplitudes of a renormalizable theory of scalar fields invariant
under a symmetry group can be made finite using only the set of counterterms
invariant to the symmetry.  This gives a complete and quite satisfactory
answer to the question posed at the beginning of Section~11.2.

The computation of the effective action in spatially varying background
fields has not been analyzed at the level of rigor involved in the proof of
the BPHZ theorem.  However, it is expected that in this situation also, the
standard set of counterterms for the symmetric theory should suffice.  We can
argue this intuitively by using the fact that the ultraviolet divergences of
Feynman diagrams are local in spacetime.  Thus, to understand the divergences
of a computation in a background $\phi_{cl}(x)$ that is smoothly varying, we
can divide spacetime into small boxes, in each of which $\phi_{cl}(x)$ is
approximately constant, and expand in the derivatives
$\partial_\mu\phi_{cl}(x)$.  In this expansion in powers of $\phi_{cl}(x)$, the
Taylor series coefficients are functional derivatives of $\Gamma$ in a
constant background, which we know can be renormalized.  The conclusion
extends to all orders, and the counterterms of the symmetric theory suffice
for the renormalization of the theory with spontaneously broken symmetry.

\subsection{Goldstone's Theorem Revisited}

As a final application of the effective action formalism, let us return to the
question of whether Goldstone's theorem is valid in the presence of quantum
corrections.  Recall that we proved this theorem at the classical level at the
end of Section~11.1: We showed in \eqref{eq:11.13} that, if the Lagrangian has
a continuous symmetry that is spontaneously broken, the matrix of second
derivatives of the classical potential $V(\phi)$ has a corresponding zero
eigenvalue.  According to Eq.~\eqref{eq:11.11}, this implies that the
classical theory contains a massless scalar particle, associated with the
spontaneously broken symmetry.

Using the effective action formalism, this argument can be repeated almost
verbatim in the full quantum field theory.  The effective potential
$V_{eff}(\phi_{cl})$ encapsulates the full solution to the theory, including
all orders of quantum corrections.  At the same time, it satisfies the general
properties of the classical potential: It is invariant to the symmetries of
the theory, and its minimum gives the vacuum expectation value of $\phi$.
This means that the argument we gave in \eqref{eq:11.13} works in exactly the
same way for $V_{eff}$ as it does for $V$: If a continuous symmetry of the
original Lagrangian is spontaneously broken by $\langle\phi\rangle$, the
matrix of second derivatives of $V_{eff}(\phi_{cl})$ has a zero eigenvalue
along the symmetry direction.

We now argue that, just as at the classical level, the presence of such a zero
eigenvalue implies the existence of a massless scalar particle.  In our
discussion of the general properties of the effective action, we showed that
its second functional derivative is the inverse propagator, and that, through
Eq.~\eqref{eq:11.97}, this derivative yields the spectrum of masses in the
quantum theory.  Let us rewrite Eq.~\eqref{eq:11.97} for a theory that
contains several scalar fields:
\begin{equation}
  \int d^4x\ e^{ip\cdot(x-y)}\ \frac{\delta^2\Gamma[\phi_{cl}]}
    {\delta\phi^i_{cl}(x)\delta\phi^j_{cl}(y)} = 0 .
  \label{eq:11.105}
\end{equation}
A particle of mass $m$ corresponds to a zero eigenvalue of this matrix
equation at $p^2 = m^2$.  Now set $p = 0$.  This implies that we
differentiate $\Gamma[\phi_{cl}]$ with respect to constant fields.  Thus, we
can replace $\Gamma[\phi_{cl}]$ by its value with constant classical fields,
which is just the effective potential.  We find that the quantum field theory
contains a scalar particle of zero mass when the matrix of second derivatives,
\begin{equation}
  \frac{\partial^2 V_{eff}}{\partial\phi^i_{cl}\partial\phi^j_{cl}} ,
\end{equation}
has a zero eigenvalue.  This completes the proof of Goldstone's theorem.

This argument for Goldstone's theorem illustrates the power of the effective
action formalism.  The formalism gives a geometrical picture of spontaneous
symmetry breaking that is valid to any order in quantum corrections.  As a
bonus, it is built up from objects that are renormalized in a simple way.
This formalism will prove useful in understanding the applications of
spontaneously broken symmetry that occur, in several different contexts,
throughout the rest of this book.

\problems

\textbf{11.1} \emph{Spin-wave theory.}
\begin{enumerate}[label=(\alph*)]
  \item Prove the following wonderful formula: Let $\phi(x)$ be a free scalar
    field with propagator $\langle T\phi(x)\phi(0)\rangle = D(x)$.  Then
    \begin{equation}
      \langle T\,e^{i\phi(x)}e^{-i\phi(0)} \rangle
      = e^{\left[ D(x) - D(0) \right]} .
    \end{equation}
    (The factor $D(0)$ gives a formally divergent adjustment of the overall
    normalization.)
  \item We can use this formula in Euclidean field theory to discuss
    correlation functions in a theory with spontaneously broken symmetry for
    $T < T_c$.  Let us consider only the simplest case of a broken $O(2)$ or
    $U(1)$ symmetry.  We can write the local spin density as a complex
    variable
    \begin{equation}
      s(x) = s_1(x) + is_2(x) .
    \end{equation}
    The global symmetry is the transformation
    \begin{equation}
      s(x) \to e^{-i\alpha}s(x) .
    \end{equation}
    If we assume that the physics freezes the modulus of $s(x)$, we can
    parametrize
    \begin{equation}
      s(x) = Ae^{i\phi(x)}
    \end{equation}
    and write an effective Lagrangian for the field $\phi(x)$.  The symmetry
    of the theory becomes the translation symmetry
    \begin{equation}
      \phi(x) \to \phi(x) - \alpha .
    \end{equation}
    Show that (for $d > 0$) the most general renormalizable Lagrangian
    consistent with this symmetry is the free field theory
    \begin{equation}
      \mathcal{L} = \tfrac{1}{2}\rho(\nabla\phi)^2 .
    \end{equation}
    In statistical mechanics, the constant $\rho$ is called the \emph{spin
    wave modulus}.  A reasonable hypothesis for $\rho$ is that it is finite
    for $T < T_c$ and tends to 0 as $T\to T_c$ from below.
  \item Compute the correlation function $\langle s(x)s^{*}(0)\rangle$.
    Adjust $A$ to give a physically sensible normalization (assuming that the
    system has a physical cutoff at the scale of one atomic spacing) and
    display the dependence of this correlation function on $x$ for
    $d = 1, 2, 3, 4$.  Explain the significance of your results.
\end{enumerate}

\textbf{11.2} \emph{A zeroth-order natural relation.}  This problem studies an
$N = 2$ linear sigma model coupled to fermions:
\begin{equation}
  \mathcal{L} = \tfrac{1}{2}(\partial_\mu\phi^i)^2 + \tfrac{1}{2}\mu^2(\phi^i)^2
    - \frac{\lambda}{4}\bigl[ (\phi^i)^2 \bigr]^2
    + \bar{\psi}(i\slashed{\partial})\psi
    - g\bar{\psi}(\phi^1 + i\gamma^5\phi^2)\psi ,
  \label{eq:11.p1}
\end{equation}
where $\phi^i$ is a two-component field, $i = 1, 2$.
\begin{enumerate}[label=(\alph*)]
  \item Show that this theory has the following global symmetry:
    \begin{equation}
      \phi^1 \to \cos\alpha\,\phi^1 - \sin\alpha\,\phi^2,\qquad
      \phi^2 \to \sin\alpha\,\phi^1 + \cos\alpha\,\phi^2,\qquad
      \psi \to e^{i\gamma^5\alpha/2}\psi .
      \label{eq:11.p2}
    \end{equation}
    Show also that the solution to the classical equations of motion with the
    minimum energy breaks this symmetry spontaneously.
  \item Denote the vacuum expectation value of the field $\phi^1$ by $v$ and
    make the change of variables
    \begin{equation}
      \phi^1(x) = (v + \sigma(x), \pi(x)) .
      \label{eq:11.p3}
    \end{equation}
    Write out the Lagrangian in these new variables, and show that the
    fermion acquires a mass given by
    \begin{equation}
      m_f = g\cdot v .
      \label{eq:11.p4}
    \end{equation}
  \item Compute the one-loop radiative correction to $m_f$, choosing
    renormalization conditions so that $v$ and $g$ (defined as the
    $\bar{\psi}\psi\pi$ vertex at zero momentum transfer) receive no radiative
    corrections.  Show that relation \eqref{eq:11.p4} receives nonzero
    corrections but that these corrections are finite.  This is in accord
    with our general discussion in Section~11.6.
\end{enumerate}

\textbf{11.3} \emph{The Gross--Neveu model.}  The Gross--Neveu model is a
model in two spacetime dimensions of fermions with a discrete chiral
symmetry:
\begin{equation}
  \mathcal{L} = \bar{\psi}_i i\slashed{\partial}\psi^i
    + \tfrac{1}{2}g^2(\bar{\psi}_i\psi^i)^2 ,
\end{equation}
with $i = 1, \ldots, N$.  The kinetic term of two-dimensional fermions is
built from matrices $\gamma^\mu$ that satisfy the two-dimensional Dirac
algebra.  These matrices can be $2\times2$:
\begin{equation}
  \gamma^0 = \sigma^2,\qquad \gamma^1 = i\sigma^1 ,
\end{equation}
where $\sigma^i$ are Pauli sigma matrices.  Define
\begin{equation}
  \gamma^5 = \gamma^0\gamma^1 = \sigma^3 ;
\end{equation}
this matrix anticommutes with the $\gamma^\mu$.
\begin{enumerate}[label=(\alph*)]
  \item Show that this theory is invariant with respect to
    \begin{equation}
      \psi^i \to \gamma^5\psi^i ,
    \end{equation}
    and that this symmetry forbids the appearance of a fermion mass.
  \item Show that this theory is renormalizable in 2 dimensions (at the level
    of dimensional analysis).
  \item Show that the functional integral for this theory can be represented
    in the following form:
    \begin{equation}
      Z = \int \mathcal{D}\sigma\ \exp\Bigl[ -\frac{i}{2}\int d^2x\ \sigma^2
        \Bigr]\ \int \mathcal{D}\bar{\psi}\,\mathcal{D}\psi\ \exp\Bigl[
        i\int d^2x\ \bar{\psi}(i\slashed{\partial} + g\sigma)\psi \Bigr] ,
    \end{equation}
    where $\sigma(x)$ (not to be confused with a Pauli matrix) is a new
    scalar field with no kinetic energy terms.
  \item Compute the leading correction to the effective potential for
    $\sigma$ by integrating over the fermion fields $\psi$.  You will
    encounter the determinant of a Dirac operator; to evaluate this
    determinant, diagonalize the operator by first going to Fourier
    components and then diagonalizing the $2\times2$ Pauli matrix associated
    with each Fourier mode.  (Alternatively, you might just take the
    determinant of this $2\times2$ matrix.)  This 1-loop contribution requires
    a renormalization proportional to $\sigma^2$ (that is, a renormalization
    of $g^2$).  Renormalize by minimal subtraction.
  \item Ignoring two-loop and higher-order contributions, minimize this
    potential.  Show that the $\sigma$ field acquires a vacuum expectation
    value which breaks the symmetry of part (a).  Convince yourself that this
    result does not depend on the particular renormalization condition
    chosen.
  \item Note that the effective potential derived in part (e) depends on $g$
    and $N$ according to the form
    \begin{equation}
      V_{eff}(\sigma_{cl}) = N\cdot f(g^2N) .
    \end{equation}
    (The overall factor of $N$ is expected in a theory with $N$ fields.)
    Construct a few of the higher-order contributions to the effective
    potential and show that they contain additional factors of $N^{-1}$ which
    suppress them if we take the limit $N\to\infty$, $(g^2N)$ fixed.  In this
    limit, the result of part (e) is unambiguous.
\end{enumerate}
